\chapter{Entropy}\label{ch:11}

\section{Introduction}\label{sec:11.1}

How unpredictable is a dynamical system? If we observe a measure-preserving transformation $T$ acting on a probability space, how much new information does each application of $T$ reveal? These are the questions that the theory of entropy answers.

The concept of entropy enters dynamical systems from two directions. From information theory, Shannon (1948) introduced entropy as a measure of uncertainty in a random variable. From the dynamical systems side, Kolmogorov (1958) and Sinai (1959) adapted Shannon's ideas to define a numerical invariant of measure-preserving transformations---the \textit{measure-theoretic entropy} (or \textit{Kolmogorov--Sinai entropy}). This invariant measures the asymptotic rate at which a dynamical system generates information.

Entropy has proven to be one of the most powerful tools in ergodic theory. It provides a way to distinguish systems that are otherwise difficult to tell apart; it connects measure-theoretic and topological dynamics through the variational principle; and it relates to Lyapunov exponents through Pesin's formula, tying together themes from Chapters 6 and 10.

We begin with Shannon's information-theoretic foundations, build up to the Kolmogorov--Sinai entropy of a measure-preserving transformation, and then explore topological entropy and the deep theorems that connect these ideas.

Throughout this chapter, $\log$ denotes the natural logarithm unless otherwise stated. Some authors use $\log_2$; the choice of base amounts to a choice of units (nats vs. bits) and does not affect any of the structural results.


\bigskip


\section{Shannon Entropy}\label{sec:11.2}

\subsection{Motivation from Information Theory}\label{sec:11.2.1}

Consider a random variable $X$ taking values in a finite set $\{x_1, x_2, \ldots, x_n\}$ with probabilities $p_1, p_2, \ldots, p_n$ respectively, where $p_i \geq 0$ and $\sum_{i=1}^n p_i = 1$. We seek a function $H(X)$ that quantifies the \textit{uncertainty} or \textit{information content} of $X$.

What properties should such a function have?

\begin{enumerate}[nosep]
  \item \textbf{Continuity}: $H$ should depend continuously on the probabilities $p_1, \ldots, p_n$.
  \item \textbf{Maximality at uniformity}: When all outcomes are equally likely ($p_i = 1/n$ for all $i$), uncertainty should be maximal.
  \item \textbf{Additivity for independent events}: If $X$ and $Y$ are independent random variables, then $H(X, Y) = H(X) + H(Y)$.
\end{enumerate}

Shannon proved that, up to a positive multiplicative constant, there is a unique function satisfying these axioms.

\begin{definition}[Shannon entropy]\label{def:11.1}
The \textit{Shannon entropy} of a discrete random variable $X$ with probability distribution $(p_1, \ldots, p_n)$ is

\[
H(X) = -\sum_{i=1}^n p_i \log p_i,
\]

where we adopt the convention $0 \log 0 = 0$ (justified by $\lim_{p \to 0^+} p \log p = 0$).

The quantity $-\log p_i$ is called the \textit{information content} (or \textit{surprisal}) of the outcome $x_i$. Rare events carry more information; certain events carry none. The entropy $H(X)$ is the expected information content:

\[
H(X) = \mathbb{E}[-\log p_X] = \sum_{i=1}^n p_i (-\log p_i).
\]
\end{definition}


\begin{example}
Consider a coin flip with $P(\text{heads}) = p$ and $P(\text{tails}) = 1 - p$. The entropy is

\[
H(p) = -p \log p - (1-p)\log(1-p).
\]

This function equals $0$ at $p = 0$ and $p = 1$ (deterministic outcomes), and achieves its maximum of $\log 2$ at $p = 1/2$ (fair coin).
\end{example}


\subsection{Basic Properties}\label{sec:11.2.2}

\begin{proposition}\label{prop:11.3}
(Properties of Shannon entropy). Let $X$ be a random variable taking $n$ values with probabilities $(p_1, \ldots, p_n)$. Then:

\textit{(i)} $H(X) \geq 0$, with equality if and only if some $p_i = 1$.

\textit{(ii)} $H(X) \leq \log n$, with equality if and only if $p_i = 1/n$ for all $i$.

\textit{(iii)} $H$ is a concave function of the probability vector $(p_1, \ldots, p_n)$.
\end{proposition}


\begin{proof}
*

\textit{(i)} Since $0 \leq p_i \leq 1$, we have $\log p_i \leq 0$, so $-p_i \log p_i \geq 0$ for each $i$. Thus $H(X) \geq 0$. We have $H(X) = 0$ if and only if $-p_i \log p_i = 0$ for every $i$, which occurs if and only if each $p_i$ is either $0$ or $1$. Since the $p_i$ sum to $1$, exactly one must equal $1$.

\textit{(ii)} We use the inequality $\log x \leq x - 1$ for $x > 0$, with equality iff $x = 1$. For each $i$ with $p_i > 0$:

\[
\log \frac{1/n}{p_i} \leq \frac{1/n}{p_i} - 1.
\]

Multiplying by $p_i$ and summing:

\[
\sum_{i=1}^n p_i \log \frac{1/n}{p_i} \leq \sum_{i=1}^n p_i \left(\frac{1/n}{p_i} - 1\right) = \sum_{i=1}^n \frac{1}{n} - \sum_{i=1}^n p_i = 1 - 1 = 0.
\]

Expanding the left side:

\[
\sum_{i=1}^n p_i \log(1/n) - \sum_{i=1}^n p_i \log p_i \leq 0,
\]

which gives $-\log n + H(X) \leq 0$, i.e., $H(X) \leq \log n$. Equality holds iff $1/(np_i) = 1$ for all $i$, i.e., $p_i = 1/n$.

\textit{(iii)} Let $\varphi(x) = -x \log x$ for $x > 0$ and $\varphi(0) = 0$. Then $\varphi''(x) = -1/x < 0$ for $x > 0$, so $\varphi$ is strictly concave. Since $H(p_1, \ldots, p_n) = \sum_{i=1}^n \varphi(p_i)$ is a sum of concave functions of the $p_i$, and each $p_i$ is a linear (hence concave) function of the probability vector, $H$ is concave.
\end{proof}


\subsection{Entropy of a Partition}\label{sec:11.2.3}

We now translate these ideas into the language of measure theory, which is the natural setting for ergodic theory.

\begin{definition}[Measurable partition]\label{def:11.4}
Let $(X, \mathcal{B}, \mu)$ be a probability space. A \textit{finite measurable partition} of $X$ is a finite collection $\alpha = \{A_1, A_2, \ldots, A_n\}$ of pairwise disjoint measurable sets with $\bigcup_{i=1}^n A_i = X$.

A partition $\alpha$ represents a finite-resolution observation of the system: we can determine which atom $A_i$ a point $x$ belongs to, but nothing finer. The entropy of $\alpha$ quantifies how much information this observation provides.
\end{definition}


\begin{definition}[Entropy of a partition]\label{def:11.5}
The \textit{entropy} of a finite measurable partition $\alpha = \{A_1, \ldots, A_n\}$ with respect to a probability measure $\mu$ is

\[
H(\alpha) = H_\mu(\alpha) = -\sum_{A \in \alpha} \mu(A) \log \mu(A).
\]

This is simply the Shannon entropy of the probability distribution $(\mu(A_1), \ldots, \mu(A_n))$.
\end{definition}


\begin{definition}[Refinement and join]\label{def:11.6}
Let $\alpha$ and $\beta$ be finite measurable partitions.

- We say $\alpha$ \textit{refines} $\beta$ (written $\alpha \succeq \beta$ or $\beta \preceq \alpha$) if every atom of $\beta$ is a union of atoms of $\alpha$.
- The \textit{join} (or \textit{common refinement}) $\alpha \vee \beta$ is the partition whose atoms are all nonempty sets of the form $A \cap B$ with $A \in \alpha$ and $B \in \beta$.

If $\alpha \succeq \beta$, then $\alpha$ is at least as fine as $\beta$: observing through $\alpha$ provides at least as much information as observing through $\beta$.
\end{definition}


\begin{proposition}\label{prop:11.7.}
If $\alpha \succeq \beta$, then $H(\alpha) \geq H(\beta)$.
\end{proposition}


\begin{proof}
* Each atom $B_j$ of $\beta$ is a disjoint union of atoms of $\alpha$: say $B_j = \bigcup_{i \in I_j} A_i$, where the index sets $I_j$ partition $\{1, \ldots, n\}$. Within each $B_j$, the concavity of $\varphi(x) = -x\log x$ gives (by Jensen's inequality applied to the identity):

\[
-\mu(B_j)\log\mu(B_j) = \varphi\!\left(\sum_{i \in I_j}\mu(A_i)\right) \leq \sum_{i \in I_j}\varphi(\mu(A_i)) = -\sum_{i \in I_j}\mu(A_i)\log\mu(A_i),
\]

where we have used the fact that $\varphi$ is concave and $\varphi$ of a sum $\geq$ the same sum when each summand is non-negative (this actually requires the \textit{superadditivity} of $\varphi$ on non-negative reals, which is the correct direction here, or one can argue via the grouping axiom for entropy). Summing over $j$:

\[
H(\beta) = \sum_j \varphi(\mu(B_j)) \leq \sum_j \sum_{i \in I_j} \varphi(\mu(A_i)) = H(\alpha). \quad \blacksquare
\]


\bigskip


\end{proof}

\section{Conditional Entropy and Joint Entropy}\label{sec:11.3}

\subsection{Conditional Entropy}\label{sec:11.3.1}

\begin{definition}[Conditional entropy]\label{def:11.8}
Let $\alpha = \{A_1, \ldots, A_m\}$ and $\beta = \{B_1, \ldots, B_n\}$ be finite measurable partitions of $(X, \mathcal{B}, \mu)$. The \textit{conditional entropy} of $\alpha$ given $\beta$ is

\[
H(\alpha \mid \beta) = -\sum_{j=1}^n \sum_{i=1}^m \mu(A_i \cap B_j) \log \frac{\mu(A_i \cap B_j)}{\mu(B_j)},
\]

where terms with $\mu(B_j) = 0$ are omitted. Equivalently:

\[
H(\alpha \mid \beta) = \sum_{j=1}^n \mu(B_j) \, H(\alpha \mid B_j),
\]

where $H(\alpha \mid B_j) = -\sum_{i=1}^m \frac{\mu(A_i \cap B_j)}{\mu(B_j)} \log \frac{\mu(A_i \cap B_j)}{\mu(B_j)}$ is the entropy of $\alpha$ computed with respect to the conditional measure $\mu(\cdot \mid B_j)$.

The conditional entropy measures how much additional information $\alpha$ provides once $\beta$ is already known.
\end{definition}


\subsection{The Chain Rule}\label{sec:11.3.2}

\begin{theorem}[Chain rule]\label{thm:11.9}
For finite measurable partitions $\alpha$ and $\beta$:

\[
H(\alpha \vee \beta) = H(\beta) + H(\alpha \mid \beta).
\]
\end{theorem}


\begin{proof}
* We compute directly:

\[
H(\alpha \vee \beta) = -\sum_{i,j} \mu(A_i \cap B_j) \log \mu(A_i \cap B_j).
\]

Writing $\mu(A_i \cap B_j) = \mu(B_j) \cdot \frac{\mu(A_i \cap B_j)}{\mu(B_j)}$ and using $\log(ab) = \log a + \log b$:

\[
H(\alpha \vee \beta) = -\sum_{i,j} \mu(A_i \cap B_j) \left[\log \mu(B_j) + \log \frac{\mu(A_i \cap B_j)}{\mu(B_j)}\right].
\]

Splitting the sum:

\[
= -\sum_{j} \log \mu(B_j) \sum_{i} \mu(A_i \cap B_j) \;-\; \sum_{i,j} \mu(A_i \cap B_j) \log \frac{\mu(A_i \cap B_j)}{\mu(B_j)}.
\]

Since $\sum_i \mu(A_i \cap B_j) = \mu(B_j)$, the first term is $-\sum_j \mu(B_j) \log \mu(B_j) = H(\beta)$. The second term is $H(\alpha \mid \beta)$ by definition.
\end{proof}


By symmetry of the join ($\alpha \vee \beta = \beta \vee \alpha$), we also have

\[
H(\alpha \vee \beta) = H(\alpha) + H(\beta \mid \alpha).
\]

\begin{corollary}[Conditioning reduces entropy]\label{cor:11.10}

\[
H(\alpha \mid \beta) \leq H(\alpha),
\]

with equality if and only if $\alpha$ and $\beta$ are independent (i.e., $\mu(A_i \cap B_j) = \mu(A_i)\mu(B_j)$ for all $i, j$).
\end{corollary}


\begin{proof}
* We have $H(\alpha \vee \beta) \leq H(\alpha) + H(\beta)$ (subadditivity, which follows from the log-sum inequality or from the non-negativity of mutual information). Combining with the chain rule:

\[
H(\beta) + H(\alpha \mid \beta) = H(\alpha \vee \beta) \leq H(\alpha) + H(\beta),
\]

so $H(\alpha \mid \beta) \leq H(\alpha)$.

For the subadditivity $H(\alpha \vee \beta) \leq H(\alpha) + H(\beta)$, one applies the inequality $\log x \leq x - 1$ as follows. Let $p_{ij} = \mu(A_i \cap B_j)$, $a_i = \mu(A_i)$, $b_j = \mu(B_j)$. Then:

\[
H(\alpha) + H(\beta) - H(\alpha \vee \beta) = \sum_{i,j} p_{ij} \log \frac{p_{ij}}{a_i b_j}.
\]

Wait---that is $-\sum_{i,j} p_{ij}\log\frac{a_i b_j}{p_{ij}}$. Applying $\log x \leq x - 1$ with $x = \frac{a_i b_j}{p_{ij}}$:

\[
-\sum_{i,j} p_{ij} \log\frac{a_i b_j}{p_{ij}} \geq -\sum_{i,j} p_{ij}\left(\frac{a_i b_j}{p_{ij}} - 1\right) = -\sum_{i,j} a_i b_j + \sum_{i,j} p_{ij} = -1 + 1 = 0.
\]

So $H(\alpha) + H(\beta) - H(\alpha \vee \beta) \geq 0$. Equality holds iff $a_i b_j / p_{ij} = 1$ for all $i, j$, i.e., $\alpha$ and $\beta$ are independent.
\end{proof}


\subsection{Further Properties}\label{sec:11.3.3}

The following properties will be needed in the construction of measure-theoretic entropy.

\begin{proposition}\label{prop:11.11.}

\textit{(i)} $H(\alpha \mid \beta) \geq 0$.

\textit{(ii)} If $\beta \preceq \gamma$ (i.e., $\gamma$ refines $\beta$), then $H(\alpha \mid \gamma) \leq H(\alpha \mid \beta)$.

\textit{(iii)} For any partitions $\alpha, \beta, \gamma$: $H(\alpha \vee \beta \mid \gamma) \leq H(\alpha \mid \gamma) + H(\beta \mid \gamma)$.

\textit{(iv)} (Chain rule, general form): $H(\alpha \vee \beta \mid \gamma) = H(\alpha \mid \gamma) + H(\beta \mid \alpha \vee \gamma)$.

Part (ii) says that conditioning on more information can only reduce (or maintain) uncertainty. Part (iii) is conditional subadditivity. We omit the proofs, which follow the same techniques as above.


\bigskip

\end{proposition}


\section{Measure-Theoretic Entropy (Kolmogorov–Sinai Entropy)}\label{sec:11.4}

We now come to the central definition of this chapter. The idea is to measure how fast a dynamical system generates information by tracking how the complexity of iterated partitions grows.

\subsection{Setup}\label{sec:11.4.1}

Let $(X, \mathcal{B}, \mu, T)$ be a measure-preserving dynamical system: $T \colon X \to X$ is measurable, $\mu$ is a probability measure, and $\mu(T^{-1}B) = \mu(B)$ for all $B \in \mathcal{B}$.

Given a finite partition $\alpha$, define its \textit{iterates} under $T$:

\[
T^{-k}\alpha = \{T^{-k}A : A \in \alpha\}.
\]

Since $T$ is measure-preserving, $\mu(T^{-k}A) = \mu(A)$, so $H(T^{-k}\alpha) = H(\alpha)$.

The join $\bigvee_{k=0}^{n-1} T^{-k}\alpha$ is the partition whose atoms are sets of the form

\[
A_{i_0} \cap T^{-1}A_{i_1} \cap \cdots \cap T^{-(n-1)}A_{i_{n-1}}, \qquad A_{i_j} \in \alpha.
\]

A point $x$ and a point $y$ lie in the same atom of this partition if and only if $x$ and $y$ visit the same sequence of atoms of $\alpha$ under the first $n$ iterates of $T$: that is, $T^k(x)$ and $T^k(y)$ belong to the same atom of $\alpha$ for $k = 0, 1, \ldots, n-1$. Thus, $\bigvee_{k=0}^{n-1} T^{-k}\alpha$ encodes the trajectory of a point up to time $n-1$, as seen through the lens of $\alpha$.

The entropy $H\!\left(\bigvee_{k=0}^{n-1} T^{-k}\alpha\right)$ measures the total uncertainty in an $n$-step trajectory. The entropy \textit{rate} is the per-step contribution:

\subsection{Definition and Existence of the Limit}\label{sec:11.4.2}

\textbf{Theorem 11.12.} Let $T$ be a measure-preserving transformation on $(X, \mathcal{B}, \mu)$ and let $\alpha$ be a finite measurable partition. Define

\[
H_n = H\!\left(\bigvee_{k=0}^{n-1} T^{-k}\alpha\right).
\]

Then:

\textit{(i)} The sequence $(H_n)$ is \textit{subadditive}: $H_{m+n} \leq H_m + H_n$ for all $m, n \geq 1$.

\textit{(ii)} The limit

\[
h_\mu(T, \alpha) = \lim_{n \to \infty} \frac{H_n}{n}
\]

exists and equals $\inf_{n \geq 1} \frac{H_n}{n}$.

\begin{proof}
*

\textit{(i)} Write $\bigvee_{k=0}^{m+n-1} T^{-k}\alpha = \bigvee_{k=0}^{m-1} T^{-k}\alpha \;\vee\; \bigvee_{k=m}^{m+n-1} T^{-k}\alpha$. By subadditivity of entropy (Corollary 11.10 and the discussion preceding it):

\[
H_{m+n} \leq H\!\left(\bigvee_{k=0}^{m-1} T^{-k}\alpha\right) + H\!\left(\bigvee_{k=m}^{m+n-1} T^{-k}\alpha\right) = H_m + H\!\left(\bigvee_{k=0}^{n-1} T^{-(k+m)}\alpha\right).
\]

Now, since $T$ is measure-preserving, the map $T^{-m}$ preserves the measure of every set, so:

\[
H\!\left(\bigvee_{k=0}^{n-1} T^{-(k+m)}\alpha\right) = H\!\left(T^{-m}\bigvee_{k=0}^{n-1} T^{-k}\alpha\right) = H\!\left(\bigvee_{k=0}^{n-1} T^{-k}\alpha\right) = H_n.
\]

Therefore $H_{m+n} \leq H_m + H_n$.

\textit{(ii)} This is a direct application of Fekete's lemma.

\end{proof}

\begin{lemma}[Fekete]\label{lem:11.13}
If $(a_n)_{n \geq 1}$ is a subadditive sequence of real numbers (i.e., $a_{m+n} \leq a_m + a_n$), then $\lim_{n \to \infty} a_n / n$ exists (in $[-\infty, \infty)$) and equals $\inf_{n \geq 1} a_n/n$.
\end{lemma}


\begin{proof}[Proof of Fekete's lemma]
of Fekete's lemma.* Fix $m \geq 1$. For any $n$, write $n = qm + r$ with $0 \leq r < m$. By subadditivity applied repeatedly, $a_n \leq q \cdot a_m + a_r$. Dividing by $n = qm + r$:

\[
\frac{a_n}{n} \leq \frac{q \cdot a_m}{qm + r} + \frac{a_r}{n}.
\]

As $n \to \infty$, $q/n \to 1/m$ and $a_r/n \to 0$ (since $0 \leq r < m$ and $a_r$ is bounded). Therefore:

\[
\limsup_{n \to \infty} \frac{a_n}{n} \leq \frac{a_m}{m}.
\]

Since this holds for every $m$:

\[
\limsup_{n \to \infty} \frac{a_n}{n} \leq \inf_{m \geq 1} \frac{a_m}{m} \leq \liminf_{n \to \infty} \frac{a_n}{n}.
\]

Therefore the limit exists and equals the infimum.
\end{proof}


Since $H_n \geq 0$ (entropy is non-negative) and $(H_n)$ is subadditive, Fekete's lemma gives $h_\mu(T, \alpha) = \lim_{n} H_n/n = \inf_n H_n/n \geq 0$. Moreover, $h_\mu(T, \alpha) \leq H_1/1 = H(\alpha)$.

\begin{remark}
There is an alternative formula:

\[
h_\mu(T, \alpha) = \lim_{n \to \infty} H\!\left(\alpha \;\Big|\; \bigvee_{k=1}^{n} T^{-k}\alpha\right).
\]

This expresses the entropy rate as the new information gained by observing the present, given the entire past. The limit is a decreasing limit (conditioning on more of the past can only decrease uncertainty), so it exists by monotonicity. We will not prove the equivalence here but refer to Walters (1982), Theorem 4.12.
\end{remark}


\subsection{The Entropy of a Transformation}\label{sec:11.4.3}

\begin{definition}[Kolmogorov–Sinai entropy]\label{def:11.15}
The \textit{measure-theoretic entropy} (or \textit{metric entropy}, or \textit{Kolmogorov--Sinai entropy}) of a measure-preserving transformation $T$ is

\[
h_\mu(T) = \sup_\alpha h_\mu(T, \alpha),
\]

where the supremum is taken over all finite measurable partitions $\alpha$ of $X$.

\textbf{Interpretation.} The entropy $h_\mu(T)$ is the maximal rate of information creation per unit time, optimized over all possible finite-resolution observations of the system. A system with $h_\mu(T) = 0$ is, in a precise sense, deterministic at the information-theoretic level: the past determines the future (up to measure zero). A system with $h_\mu(T) > 0$ is genuinely unpredictable---new information is created at every step.
\end{definition}


\subsection{Basic Properties of Metric Entropy}\label{sec:11.4.4}

\begin{proposition}\label{prop:11.16.}

\textit{(i)} $h_\mu(T) \geq 0$.

\textit{(ii)} $h_\mu(T^n) = n \cdot h_\mu(T)$ for all $n \geq 1$.

\textit{(iii)} If $T$ is invertible, $h_\mu(T^{-1}) = h_\mu(T)$.

\textit{(iv)} If $(X_1, \mu_1, T_1)$ and $(X_2, \mu_2, T_2)$ are measure-theoretically isomorphic, then $h_{\mu_1}(T_1) = h_{\mu_2}(T_2)$.
\end{proposition}


\begin{proof}[Proof sketch]
sketch.* (i) is immediate. For (ii), observe that $\bigvee_{k=0}^{n-1}(T^n)^{-k}\alpha = \bigvee_{k=0}^{n-1} T^{-nk}\alpha$, and a careful computation shows $h_\mu(T^n, \alpha) = h_\mu(T, \bigvee_{k=0}^{n-1} T^{-k}\alpha)$; taking the supremum over $\alpha$ yields the result. Part (iii) follows from (ii) and the observation that $h_\mu(\mathrm{id}) = 0$. Part (iv) follows because an isomorphism establishes a bijection between partitions preserving all relevant quantities. See Walters (1982), Chapter 4, for complete proofs. $\blacksquare$

Property (iv) is the reason entropy is so important: it is a \textit{conjugacy invariant}. If two systems have different entropies, they cannot be isomorphic.


\bigskip


\end{proof}

\section{Computing Entropy: The Generator Theorem}\label{sec:11.5}

The supremum in the definition of $h_\mu(T)$ is taken over all finite partitions---an enormous collection. The Kolmogorov--Sinai theorem tells us that this supremum is achieved by any \textit{generating partition}, reducing the computation of entropy to a single partition.

\subsection{Generating Partitions}\label{sec:11.5.1}

\begin{definition}[Generating partition]\label{def:11.17}
A finite measurable partition $\alpha$ is a \textit{generator} (or \textit{generating partition}) for a measure-preserving transformation $T$ if the $\sigma$-algebra generated by the iterated partitions $\{T^{-k}\alpha : k \geq 0\}$ is the full $\sigma$-algebra $\mathcal{B}$ (up to sets of measure zero). That is:

\[
\bigvee_{k=0}^{\infty} T^{-k}\alpha \doteq \mathcal{B},
\]

where $\doteq$ means equality modulo $\mu$-null sets. If $T$ is invertible, we may also require the \textit{two-sided} generating condition:

\[
\bigvee_{k=-\infty}^{\infty} T^{-k}\alpha \doteq \mathcal{B}.
\]

Informally, $\alpha$ is a generator if, by observing the trajectory of a point through the atoms of $\alpha$ for all time, one can (in principle) determine the point exactly.
\end{definition}


\begin{example}
For the full shift $\sigma$ on $\Sigma_k = \{1, 2, \ldots, k\}^{\mathbb{Z}}$, the partition $\alpha = \{[i] : i = 1, \ldots, k\}$, where $[i] = \{x \in \Sigma_k : x_0 = i\}$ is the cylinder set at position $0$, is a generating partition. Indeed, knowing $x_0, x_1, x_2, \ldots$ (by applying $\sigma^{-k}\alpha$ for all $k \geq 0$) determines the entire forward sequence; for the two-sided shift, knowing $x_n$ for all $n \in \mathbb{Z}$ determines $x$ exactly.
\end{example}


\begin{example}
For the doubling map $T(x) = 2x \pmod{1}$ on $[0,1)$ with Lebesgue measure, the partition $\alpha = \{[0, 1/2), [1/2, 1)\}$ is a generator. The atom containing $T^k(x)$ tells us the $(k+1)$-th digit in the binary expansion of $x$; knowing all digits determines $x$ (up to the countable set of dyadic rationals, which has measure zero).
\end{example}


\subsection{The Kolmogorov–Sinai Theorem}\label{sec:11.5.2}

\begin{theorem}[Kolmogorov–Sinai]\label{thm:11.20}
Let $T$ be a measure-preserving transformation on $(X, \mathcal{B}, \mu)$ and let $\alpha$ be a generating partition (one-sided for non-invertible $T$, two-sided for invertible $T$). Then:

\[
h_\mu(T) = h_\mu(T, \alpha).
\]

This theorem is the workhorse for computing entropy: instead of taking a supremum over all partitions, we need only find a single generator and compute $h_\mu(T, \alpha)$.
\end{theorem}


\begin{proof}[Proof sketch]
sketch.* The key idea is that for any finite partition $\beta$, if $\alpha$ is a generator, then the partitions $\bigvee_{k=0}^{n-1} T^{-k}\alpha$ eventually approximate $\beta$ arbitrarily well (in the sense that $\beta$ can be approximated by a partition measurable with respect to $\bigvee_{k=0}^{n-1} T^{-k}\alpha$). One shows that for any $\epsilon > 0$ and any finite partition $\beta$, if $n$ is large enough, then $h_\mu(T, \beta) \leq h_\mu(T, \alpha) + \epsilon$. Taking $\epsilon \to 0$ and the supremum over $\beta$ gives $h_\mu(T) \leq h_\mu(T, \alpha)$. The reverse inequality $h_\mu(T) \geq h_\mu(T, \alpha)$ holds by definition. See Petersen (1983), Theorem 4.17, or Walters (1982), Theorem 4.18, for the complete proof. $\blacksquare$


\bigskip


\end{proof}

\section{Examples}\label{sec:11.6}

We now compute entropy explicitly for several important dynamical systems.

\subsection{Bernoulli Shifts}\label{sec:11.6.1}

\textbf{The fair coin: $B(1/2, 1/2)$.} Consider the Bernoulli shift $\sigma$ on $\{0, 1\}^{\mathbb{Z}}$ (or $\{0,1\}^{\mathbb{N}}$) where each coordinate independently takes the value $0$ or $1$ with equal probability $1/2$. The partition $\alpha = \{[0], [1]\}$ into cylinder sets at position $0$ is a generator (Example 11.18).

We compute $h_\mu(\sigma, \alpha)$ directly. The join $\bigvee_{k=0}^{n-1} \sigma^{-k}\alpha$ partitions the space according to the first $n$ coordinates $(x_0, x_1, \ldots, x_{n-1})$. There are $2^n$ atoms, each of measure $(1/2)^n = 2^{-n}$. Therefore:

\[
H\!\left(\bigvee_{k=0}^{n-1} \sigma^{-k}\alpha\right) = -\sum_{w \in \{0,1\}^n} 2^{-n} \log 2^{-n} = -2^n \cdot 2^{-n} \cdot (-n \log 2) = n \log 2.
\]

Thus:

\[
h_\mu(\sigma) = h_\mu(\sigma, \alpha) = \lim_{n \to \infty} \frac{n \log 2}{n} = \log 2.
\]

\textbf{General Bernoulli shift $B(p_1, \ldots, p_k)$.} Now consider the Bernoulli shift on $\{1, 2, \ldots, k\}^{\mathbb{Z}}$ where the $i$-th symbol appears with probability $p_i$. Again, $\alpha = \{[1], [2], \ldots, [k]\}$ is a generating partition. Each atom of $\bigvee_{k=0}^{n-1}\sigma^{-k}\alpha$ corresponds to a word $w = (w_0, w_1, \ldots, w_{n-1}) \in \{1, \ldots, k\}^n$ with measure $p_{w_0} p_{w_1} \cdots p_{w_{n-1}}$.

\[
H\!\left(\bigvee_{j=0}^{n-1}\sigma^{-j}\alpha\right) = -\sum_{w \in \{1,\ldots,k\}^n} \left(\prod_{j=0}^{n-1} p_{w_j}\right) \log\!\left(\prod_{j=0}^{n-1} p_{w_j}\right).
\]

Since the coordinates are independent:

\[
= -\sum_{w} \prod_{j} p_{w_j} \sum_{j=0}^{n-1} \log p_{w_j} = -\sum_{j=0}^{n-1} \sum_{w} \prod_{\ell} p_{w_\ell} \cdot \log p_{w_j}.
\]

For each fixed $j$, summing over $w$ with $w_j = i$ gives $p_i \log p_i$ (the other coordinates sum to $1$ independently). So:

\[
= -\sum_{j=0}^{n-1} \sum_{i=1}^k p_i \log p_i = -n \sum_{i=1}^k p_i \log p_i = n \cdot H(\alpha).
\]

Therefore:

\[
\boxed{h_\mu(\sigma) = -\sum_{i=1}^k p_i \log p_i.}
\]

The entropy of a Bernoulli shift is exactly the Shannon entropy of its weight vector.

\subsection{Irrational Rotations}\label{sec:11.6.2}

\begin{proposition}\label{prop:11.21.}
Let $R_\theta \colon \mathbb{T} \to \mathbb{T}$ be the rotation $x \mapsto x + \theta \pmod{1}$ on the circle $\mathbb{T} = \mathbb{R}/\mathbb{Z}$, equipped with Lebesgue measure $\lambda$. Then $h_\lambda(R_\theta) = 0$.
\end{proposition}


\begin{proof}
* Let $\alpha$ be any finite partition of $\mathbb{T}$ into intervals (or any finite measurable partition). We will show $h_\lambda(R_\theta, \alpha) = 0$. It suffices to show this for interval partitions, since every finite partition can be approximated by one.

Let $\alpha = \{A_1, \ldots, A_m\}$ where each $A_i$ is an interval (or a finite union of intervals). The partition $\alpha_n = \bigvee_{k=0}^{n-1} R_\theta^{-k}\alpha$ has at most $mn$ atoms (each application of $R_\theta^{-1}$ introduces at most $m$ new boundary points, so after $n$ steps there are at most $mn$ intervals). In general, for a partition into intervals, $|\alpha_n| \leq mn$.

By the bound $H(\alpha_n) \leq \log|\alpha_n|$:

\[
h_\lambda(R_\theta, \alpha) = \lim_{n \to \infty} \frac{H(\alpha_n)}{n} \leq \lim_{n \to \infty} \frac{\log(mn)}{n} = 0.
\]

Since $\alpha$ was arbitrary, $h_\lambda(R_\theta) = \sup_\alpha h_\lambda(R_\theta, \alpha) = 0$.
\end{proof}


\begin{remark}
More generally, \textit{every isometry} of a compact metric space has zero topological entropy (and hence zero metric entropy for every invariant measure). Zero entropy does not mean the system is uninteresting---irrational rotations are uniquely ergodic and exhibit equidistribution---but it does mean that the system is, in an information-theoretic sense, deterministic: the past determines the future.
\end{remark}


\subsection{Arnold's Cat Map}\label{sec:11.6.3}

The \textit{Arnold cat map} is the automorphism $T_A \colon \mathbb{T}^2 \to \mathbb{T}^2$ induced by the matrix

\[
A = \begin{pmatrix} 2 & 1 \\ 1 & 1 \end{pmatrix}
\]

acting on $\mathbb{R}^2/\mathbb{Z}^2$, preserving Lebesgue measure $\lambda$ on $\mathbb{T}^2$.

The eigenvalues of $A$ are

\[
\lambda_{1,2} = \frac{3 \pm \sqrt{5}}{2}.
\]

Thus $\lambda_1 = \frac{3 + \sqrt{5}}{2} > 1$ and $\lambda_2 = \frac{3 - \sqrt{5}}{2} = 1/\lambda_1 \in (0,1)$. The Lyapunov exponents are $\log \lambda_1 > 0$ and $\log \lambda_2 < 0$.

\begin{proposition}\label{prop:11.23.}
The entropy of Arnold's cat map with respect to Lebesgue measure is

\[
h_\lambda(T_A) = \log\!\left(\frac{3 + \sqrt{5}}{2}\right).
\]
\end{proposition}


\begin{proof}
* This can be computed directly using the theory of toral automorphisms. More generally, for a hyperbolic toral automorphism induced by a matrix $A \in \mathrm{GL}(d, \mathbb{Z})$, the entropy with respect to Lebesgue measure is

\[
h_\lambda(T_A) = \sum_{|\lambda_i| > 1} \log |\lambda_i|,
\]

where the sum is over eigenvalues $\lambda_i$ of $A$ with absolute value greater than $1$ (counted with multiplicity).

One approach is via Pesin's formula (Section 11.9), which applies because Lebesgue measure is the SRB measure for toral automorphisms. The positive Lyapunov exponent is $\log \lambda_1 = \log\frac{3+\sqrt{5}}{2}$, and Pesin's formula gives

\[
h_\lambda(T_A) = \log\frac{3 + \sqrt{5}}{2}.
\]

Alternatively, one can construct a Markov partition for $T_A$ and use the Kolmogorov--Sinai theorem. See Katok and Hasselblatt (1995), Section 3.2, for details.
\end{proof}


Numerically, $h_\lambda(T_A) = \log\!\left(\frac{3+\sqrt{5}}{2}\right) \approx 0.9624$.

\subsection{The Doubling Map}\label{sec:11.6.4}

The \textit{doubling map} is $T \colon [0,1) \to [0,1)$ defined by $T(x) = 2x \pmod{1}$, preserving Lebesgue measure.

\begin{proposition}\label{prop:11.24.}
$h_\lambda(T) = \log 2$.
\end{proposition}


\begin{proof}
* The partition $\alpha = \{[0, 1/2), [1/2, 1)\}$ is a generator (Example 11.19). The join $\bigvee_{k=0}^{n-1} T^{-k}\alpha$ partitions $[0,1)$ into $2^n$ equal intervals of the form $[j/2^n, (j+1)/2^n)$, $j = 0, 1, \ldots, 2^n - 1$. Each has Lebesgue measure $2^{-n}$. Therefore:

\[
H\!\left(\bigvee_{k=0}^{n-1} T^{-k}\alpha\right) = -2^n \cdot 2^{-n} \log 2^{-n} = n \log 2.
\]

So $h_\lambda(T, \alpha) = \log 2$, and by the Kolmogorov--Sinai theorem, $h_\lambda(T) = \log 2$.
\end{proof}


\begin{remark}
The doubling map is measurably isomorphic to the one-sided Bernoulli shift $B(1/2, 1/2)$. The isomorphism is given by the binary expansion map: $x \mapsto (x_1, x_2, x_3, \ldots)$ where $x = \sum x_i 2^{-i}$. That both systems have entropy $\log 2$ is consistent with the fact that entropy is a conjugacy invariant.
\end{remark}


\subsection{Entropy as a Conjugacy Invariant}\label{sec:11.6.5}

\textbf{Theorem 11.26.} Measure-theoretic entropy is an invariant of measure-theoretic isomorphism. That is, if $\varphi \colon (X_1, \mu_1) \to (X_2, \mu_2)$ is a measure-preserving isomorphism with $\varphi \circ T_1 = T_2 \circ \varphi$ (a.e.), then $h_{\mu_1}(T_1) = h_{\mu_2}(T_2)$.

This is Proposition 11.16(iv). As a consequence:

\textbf{Corollary 11.27.} The Bernoulli shifts $B(1/2, 1/2)$ and $B(1/3, 1/3, 1/3)$ are \textit{not} isomorphic.

\begin{proof}
* $h(B(1/2, 1/2)) = \log 2 \approx 0.693$ and $h(B(1/3, 1/3, 1/3)) = \log 3 \approx 1.099$. Since $\log 2 \neq \log 3$, the systems cannot be isomorphic. $\blacksquare$

Before entropy was available, distinguishing these two systems was an open problem. Both are ergodic, mixing, and have similar qualitative behavior. Entropy resolved the question definitively.


\bigskip


\end{proof}

\section{Topological Entropy}\label{sec:11.7}

We now turn to an entropy concept that lives in the topological category rather than the measure-theoretic one. Topological entropy, introduced by Adler, Konheim, and McAndrew (1965) and reformulated by Bowen (1971) and Dinaburg (1970), measures the complexity of a continuous map without reference to any invariant measure.

\subsection{Definition via Separated Sets}\label{sec:11.7.1}

Let $(X, d)$ be a compact metric space and $T \colon X \to X$ a continuous map.

\begin{definition}[Bowen metric]\label{def:11.28}
For $n \geq 1$, the \textit{Bowen metric} (or \textit{dynamical metric}) is

\[
d_n(x, y) = \max_{0 \leq k \leq n-1} d(T^k x, T^k y).
\]

Two points are close in the Bowen metric $d_n$ if their first $n$ iterates remain close.
\end{definition}


\begin{definition}[Separated and spanning sets]\label{def:11.29}
Let $\epsilon > 0$ and $n \geq 1$.

- A set $E \subset X$ is *$(n, \epsilon)$-separated* if for every pair $x \neq y$ in $E$, $d_n(x, y) > \epsilon$: the orbits of any two distinct points in $E$ differ by more than $\epsilon$ within the first $n$ steps.
- A set $F \subset X$ is *$(n, \epsilon)$-spanning* if for every $x \in X$, there exists $y \in F$ with $d_n(x, y) \leq \epsilon$: the orbits of points in $F$ approximate every orbit to within $\epsilon$.

Let $s_n(\epsilon, T)$ denote the maximal cardinality of an $(n, \epsilon)$-separated set, and $r_n(\epsilon, T)$ the minimal cardinality of an $(n, \epsilon)$-spanning set. Since $X$ is compact, both are finite.
\end{definition}


\textbf{Lemma 11.30.} $r_n(\epsilon, T) \leq s_n(\epsilon, T) \leq r_n(\epsilon/2, T)$.

\begin{proof}
* For the first inequality: a maximal $(n,\epsilon)$-separated set is automatically $(n,\epsilon)$-spanning (if some $x$ were not within $d_n$-distance $\epsilon$ of any point in $E$, we could add $x$ to $E$, contradicting maximality). For the second: if $F$ is $(n, \epsilon/2)$-spanning, then any $(n, \epsilon)$-separated set $E$ has $|E| \leq |F|$, because the map assigning to each $e \in E$ a nearest point in $F$ (within $d_n$-distance $\epsilon/2$) must be injective. $\blacksquare$

\end{proof}

\begin{definition}[Topological entropy]\label{def:11.31}
The \textit{topological entropy} of $T$ is

\[
h_{\mathrm{top}}(T) = \lim_{\epsilon \to 0} \limsup_{n \to \infty} \frac{1}{n} \log s_n(\epsilon, T).
\]

By Lemma 11.30, this equals $\lim_{\epsilon \to 0} \limsup_{n \to \infty} \frac{1}{n}\log r_n(\epsilon, T)$.
\end{definition}


\begin{remark}
As $\epsilon$ decreases, $s_n(\epsilon, T)$ can only increase, so $\frac{1}{n}\log s_n(\epsilon, T)$ is non-decreasing in $1/\epsilon$. Therefore the outer limit $\lim_{\epsilon \to 0}$ exists (as a limit in $[0, \infty]$). In many cases the $\limsup$ is actually a limit.
\end{remark}


\subsection{Definition via Open Covers}\label{sec:11.7.2}

There is an equivalent definition using open covers, which was the original approach. Let $\mathcal{U}$ be a finite open cover of $X$. Define $N(\mathcal{U})$ to be the minimal cardinality of a subcover, and for two covers $\mathcal{U}, \mathcal{V}$, define the join $\mathcal{U} \vee \mathcal{V} = \{U \cap V : U \in \mathcal{U}, V \in \mathcal{V}\}$.

\[
h_{\mathrm{top}}(T) = \sup_{\mathcal{U}} \lim_{n \to \infty} \frac{1}{n} \log N\!\left(\bigvee_{k=0}^{n-1} T^{-k}\mathcal{U}\right).
\]

The limit exists by a subadditivity argument analogous to that in Theorem 11.12.

\subsection{Examples}\label{sec:11.7.3}

\begin{example}
(Full shift on $k$ symbols). Let $\Sigma_k = \{1, \ldots, k\}^{\mathbb{Z}}$ with the product topology (metrized by, e.g., $d(x,y) = 2^{-\min\{|n| : x_n \neq y_n\}}$) and let $\sigma$ be the shift map. For any $\epsilon < 1$ small enough that the $0$-th coordinate is determined, an $(n, \epsilon)$-separated set can distinguish the coordinates $0, 1, \ldots, n-1$, so $s_n(\epsilon, \sigma) = k^n$. Therefore:

\[
h_{\mathrm{top}}(\sigma) = \lim_{n \to \infty} \frac{1}{n} \log k^n = \log k.
\]
\end{example}


\begin{example}
(Topological entropy of rotations). For an irrational rotation $R_\theta$ on the circle, $R_\theta$ is an isometry: $d(R_\theta x, R_\theta y) = d(x,y)$. Therefore $d_n(x,y) = d(x,y)$ for all $n$, and $s_n(\epsilon, R_\theta) = s_1(\epsilon, R_\theta)$ is independent of $n$. Thus:

\[
h_{\mathrm{top}}(R_\theta) = \lim_{n \to \infty} \frac{1}{n}\log s_1(\epsilon, R_\theta) = 0.
\]

This is consistent with the measure-theoretic result (Proposition 11.21).
\end{example}


\begin{example}
(Doubling map). The doubling map $T(x) = 2x \pmod 1$ on $[0,1)$ has topological entropy $h_{\mathrm{top}}(T) = \log 2$. To see this, note that $T$ is semiconjugate to the one-sided shift on $\{0,1\}^{\mathbb{N}}$ via binary expansion, and the shift on two symbols has topological entropy $\log 2$.


\bigskip

\end{example}


\section{The Variational Principle}\label{sec:11.8}

The variational principle is the bridge between measure-theoretic and topological entropy.

\begin{theorem}[Variational principle]\label{thm:11.36}
Let $T \colon X \to X$ be a continuous map on a compact metrizable space. Then:

\[
h_{\mathrm{top}}(T) = \sup_{\mu \in \mathcal{M}_T(X)} h_\mu(T),
\]

where $\mathcal{M}_T(X)$ denotes the set of all $T$-invariant Borel probability measures on $X$.

This remarkable theorem, proved independently by Goodwyn (1969, the $\geq$ direction---actually Goodwyn proved $h_\mu(T) \leq h_{\mathrm{top}}(T)$ for all $\mu$), Dinaburg (1970), and Goodman (1971, the $\leq$ direction), says that topological entropy is the supremum of metric entropies over all invariant measures.
\end{theorem}


\begin{proof}[Proof outline]
outline.* We sketch the two directions.

\textbf{Direction 1}: $h_\mu(T) \leq h_{\mathrm{top}}(T)$ for every $\mu \in \mathcal{M}_T(X)$.

Let $\alpha$ be a finite partition and $\epsilon > 0$ such that every atom of $\alpha$ has diameter less than $\epsilon$ (possible by refining $\alpha$). Then the atoms of $\bigvee_{k=0}^{n-1} T^{-k}\alpha$ are contained in $d_n$-balls of radius $\epsilon$. An $(n, \epsilon)$-separated set $E$ can intersect each atom of $\bigvee_{k=0}^{n-1} T^{-k}\alpha$ in at most one point (since any two points in the same atom have $d_n$-distance $\leq \epsilon$, but separated points have $d_n$-distance $> \epsilon$). Wait---actually this requires $\epsilon$ to be the diameter of atoms. Let us be more precise.

Take $\alpha$ such that every atom of $\alpha$ has diameter less than $\epsilon$. Then the atoms of $\alpha_n = \bigvee_{k=0}^{n-1} T^{-k}\alpha$ have diameter less than $\epsilon$ in the $d_n$ metric. The number of atoms of $\alpha_n$ with positive measure is at most $|\alpha_n|$, and $|\alpha_n|$ is bounded below by $s_n(\epsilon, T)$... actually, the correct argument uses the fact that $H(\alpha_n) \leq \log |\alpha_n|$ and that $|\alpha_n|$ is related to spanning/separated sets. More carefully: the number of nonempty atoms of $\alpha_n$ is at least $r_n(\epsilon, T)$ (since these atoms form an $(n,\epsilon)$-spanning cover in some sense). The complete argument requires the Misiurewicz proof; see Walters (1982), Theorem 8.6.

\textbf{Direction 2}: $h_{\mathrm{top}}(T) \leq \sup_\mu h_\mu(T)$.

This is the deeper direction. For each $n$ and $\epsilon$, let $E_n$ be a maximal $(n, \epsilon)$-separated set. Define an atomic measure $\nu_n = \frac{1}{|E_n|}\sum_{x \in E_n}\delta_x$, and let $\mu_n = \frac{1}{n}\sum_{k=0}^{n-1}(T^k)_*\nu_n$ be its time average. Take a weak-* limit point $\mu$; it is $T$-invariant. A careful analysis shows $h_\mu(T) \geq \limsup_{n}\frac{1}{n}\log|E_n| - \delta(\epsilon)$ where $\delta(\epsilon) \to 0$ as $\epsilon \to 0$. See Walters (1982), Theorem 8.6, for the full proof.
\end{proof}


\begin{definition}[Measure of maximal entropy]\label{def:11.37}
An invariant measure $\mu$ is called a \textit{measure of maximal entropy} (or \textit{MME}) if $h_\mu(T) = h_{\mathrm{top}}(T)$. Such a measure always exists when $X$ is compact and $T$ is continuous (by upper semicontinuity of the entropy map and compactness of $\mathcal{M}_T(X)$ in the weak-* topology).
\end{definition}


\begin{example}
For the full shift on $\{1, \ldots, k\}^{\mathbb{Z}}$, the measure of maximal entropy is the uniform Bernoulli measure $(1/k, \ldots, 1/k)$. Indeed, $h_{\mathrm{top}}(\sigma) = \log k$ and the entropy of $B(1/k, \ldots, 1/k)$ is $-k \cdot \frac{1}{k}\log\frac{1}{k} = \log k$.


\bigskip

\end{example}


\section{Pesin's Formula (The Entropy Formula)}\label{sec:11.9}

Pesin's formula connects entropy to Lyapunov exponents, providing a bridge between the information-theoretic and geometric descriptions of chaos. Recall from Chapter 6 that the Lyapunov exponents $\lambda_1(x) \geq \lambda_2(x) \geq \cdots \geq \lambda_d(x)$ measure the asymptotic rates of expansion and contraction along different directions.

\subsection{The Ruelle Inequality}\label{sec:11.9.1}

\begin{theorem}[Ruelle inequality, 1978]\label{thm:11.39}
Let $M$ be a compact smooth Riemannian manifold and $T \colon M \to M$ a $C^1$ diffeomorphism preserving a Borel probability measure $\mu$. Then:

\[
h_\mu(T) \leq \int_M \sum_{\lambda_i(x) > 0} \lambda_i(x) \, d\mu(x),
\]

where $\lambda_1(x) \geq \cdots \geq \lambda_d(x)$ are the Lyapunov exponents at $x$ (counted with multiplicity).

The Ruelle inequality holds for \textit{all} $T$-invariant probability measures $\mu$. The right side is the average exponential rate at which $T$ stretches volume in the unstable directions---an upper bound on how fast information can be created.
\end{theorem}


\begin{proof}[Proof idea]
idea.* The key is that entropy measures distinguishability of orbits, and Lyapunov exponents control how fast nearby orbits separate. A ball of radius $\epsilon$ in the Bowen metric $d_n$ has volume roughly $\epsilon^d \cdot \prod_i e^{-n\lambda_i^+(x)}$ near $x$ (where $\lambda_i^+$ denotes the positive part). The number of such balls needed to cover $M$ is roughly $\prod_i e^{n\lambda_i^+(x)}$, giving the bound. The rigorous proof uses the theory of Lyapunov charts and Pesin's theory of non-uniform hyperbolicity. See Katok and Hasselblatt (1995), Supplement S.3.

\end{proof}

\subsection{Pesin's Formula}\label{sec:11.9.2}

\begin{theorem}[Pesin's entropy formula]\label{thm:11.40}
Let $M$ be a compact smooth Riemannian manifold and $T \colon M \to M$ a $C^{1+\alpha}$ diffeomorphism (i.e., $C^1$ with $\alpha$-H\"older continuous derivative for some $\alpha > 0$). If $\mu$ is an \textit{SRB measure} (Sinai--Ruelle--Bowen measure)---in particular, if $\mu$ is absolutely continuous with respect to the Riemannian volume on the unstable manifolds---then equality holds in the Ruelle inequality:

\[
h_\mu(T) = \int_M \sum_{\lambda_i(x) > 0} \lambda_i(x) \, d\mu(x).
\]

This was proved by Pesin (1977) for the case when $\mu$ is equivalent to the Riemannian volume (Lebesgue measure on the manifold), and extended by Ledrappier and Strelcyn (1982) and Ledrappier and Young (1985) to the SRB case.
\end{theorem}


\begin{remark}
The converse direction was established by Ledrappier and Young (1985): if $h_\mu(T) = \int \sum_{\lambda_i > 0} \lambda_i \, d\mu$, then $\mu$ has absolutely continuous conditional measures on unstable manifolds---i.e., $\mu$ is an SRB measure.
\end{remark}


\subsection{Examples Revisited}\label{sec:11.9.3}

\textbf{Arnold's cat map.} For $T_A$ with $A = \bigl(\begin{smallmatrix} 2 & 1 \\ 1 & 1\end{smallmatrix}\bigr)$, the Lyapunov exponents with respect to Lebesgue measure are constant: $\lambda_1 = \log\frac{3+\sqrt{5}}{2} > 0$ and $\lambda_2 = \log\frac{3-\sqrt{5}}{2} < 0$. Lebesgue measure is the SRB measure (it is smooth, hence absolutely continuous). Pesin's formula gives:

\[
h_\lambda(T_A) = \lambda_1 = \log\frac{3+\sqrt{5}}{2}.
\]

This confirms the result of Proposition 11.23.

\textbf{Doubling map.} The doubling map $T(x) = 2x \pmod 1$ is not invertible, but there is a version of Pesin's formula for endomorphisms. The single Lyapunov exponent is $\lambda = \log 2$ (since $T'(x) = 2$ everywhere). Lebesgue measure is the natural measure, and:

\[
h_\lambda(T) = \log 2,
\]

consistent with Proposition 11.24.

\subsection{Significance}\label{sec:11.9.4}

Pesin's formula is a deep result that ties together three fundamental quantities:

\begin{enumerate}[nosep]
  \item \textbf{Entropy} (information theory / ergodic theory),
  \item \textbf{Lyapunov exponents} (dynamical instability / chaos),
  \item \textbf{SRB measures} (physically observable measures).
\end{enumerate}

It tells us that for "physically relevant" measures, the entropy is entirely accounted for by the expansion rates in the unstable directions. No information is wasted, so to speak: the Ruelle inequality becomes an equality. Systems where entropy is strictly less than $\int \sum \lambda_i^+ \, d\mu$ are, in a sense, not in equilibrium---the measure $\mu$ fails to "see" all the complexity that the dynamics is capable of producing.


\bigskip


\section{Ornstein's Theorem}\label{sec:11.10}

We conclude with one of the deepest and most celebrated results in ergodic theory.

\subsection{The Classification Problem}\label{sec:11.10.1}

After Kolmogorov introduced entropy in 1958, the immediate question was: \textit{Is entropy a complete invariant for Bernoulli shifts?} That is, if two Bernoulli shifts have the same entropy, must they be isomorphic?

Entropy can \textit{distinguish} non-isomorphic Bernoulli shifts (Corollary 11.27), but it was not clear whether it was a strong enough invariant to \textit{classify} them. For example, $B(1/2, 1/2)$ and $B(1/4, 1/4, 1/4, 1/4)$ both have entropy $\log 2$. Are they isomorphic?

\subsection{The Theorem}\label{sec:11.10.2}

\begin{theorem}[Ornstein, 1970]\label{thm:11.42}
Two Bernoulli shifts are measurably isomorphic if and only if they have the same entropy.

In other words, entropy is a \textit{complete invariant} for the class of Bernoulli shifts: it classifies them up to measure-theoretic isomorphism.
\end{theorem}


\textbf{Corollary 11.43.} $B(1/2, 1/2) \cong B(1/4, 1/4, 1/4, 1/4)$, since both have entropy $\log 2$.

This is a remarkable and non-obvious statement. These two shifts have different alphabet sizes (2 vs. 4) and different probability vectors, yet they are measure-theoretically the same system. The isomorphism is highly non-trivial and cannot be described by any simple formula.

\subsection{Significance and Methods}\label{sec:11.10.3}

The proof of Ornstein's theorem introduced a powerful new technique: the theory of *$\bar{d}$-distance\textit{ (or }finitary coding\textit{) between processes. Ornstein defined a metric on the space of stationary processes and showed that Bernoulli shifts with the same entropy can be approximated and matched via this metric. The key technical tool is the }Ornstein copying lemma*.

Ornstein's work opened a rich program of classifying dynamical systems. Several important classes of systems were subsequently shown to be Bernoulli (and hence classified by entropy):

\begin{itemize}[nosep]
  \item Geodesic flows on negatively curved manifolds (Ornstein and Weiss, 1973).
  \item Hyperbolic toral automorphisms (Katznelson, 1971).
  \item Mixing Markov chains (Friedman and Ornstein, 1970).
  \item The Gauss map $x \mapsto \{1/x\}$ on $(0,1]$ with the Gauss measure (a result that follows from its being a natural extension of a one-sided Bernoulli process).
\end{itemize}

The full proof of Ornstein's theorem is beyond our scope (it requires substantial machinery from the theory of joinings and the $\bar{d}$-metric), but we refer the interested reader to Ornstein (1974), Shields (1973), or the excellent exposition in Petersen (1983), Chapter 6.


\bigskip


\section{Summary}\label{sec:11.11}

| System | Entropy |
|--------|---------|
| Bernoulli shift $B(p_1, \ldots, p_k)$ | $-\sum p_i \log p_i$ |
| Irrational rotation $R_\theta$ | $0$ |
| Arnold's cat map | $\log\frac{3+\sqrt{5}}{2}$ |
| Doubling map $x \mapsto 2x \pmod 1$ | $\log 2$ |
| Full shift on $k$ symbols (topological) | $\log k$ |

The theory of entropy provides a hierarchy of results:

\begin{enumerate}[nosep]
  \item \textbf{Shannon entropy} quantifies the uncertainty in a single observation (partition).
  \item \textbf{Kolmogorov--Sinai entropy} measures the rate of information production of a measure-preserving system.
  \item \textbf{Topological entropy} measures the complexity of a continuous map without reference to a measure.
  \item \textbf{The variational principle} connects (2) and (3): $h_{\mathrm{top}}(T) = \sup_\mu h_\mu(T)$.
  \item \textbf{Pesin's formula} connects entropy to Lyapunov exponents for SRB measures.
  \item \textbf{Ornstein's theorem} shows that entropy completely classifies Bernoulli shifts.


\bigskip

\end{enumerate}

\section*{Exercises}

\begin{exercise}
Compute the Shannon entropy $H(p, 1-p)$ for $p \in [0,1]$ and verify that it is maximized at $p = 1/2$ with value $\log 2$. Sketch the graph.
\end{exercise}


\begin{exercise}
Let $\alpha = \{A_1, A_2, A_3\}$ be a partition of a probability space with $\mu(A_1) = 1/2$, $\mu(A_2) = 1/3$, $\mu(A_3) = 1/6$. Compute $H(\alpha)$. Compare this to $\log 3$ and explain why $H(\alpha) < \log 3$.
\end{exercise}


\begin{exercise}
(Subadditivity in action). Let $T$ be a measure-preserving transformation and $\alpha$ a finite partition. Define $H_n = H\!\left(\bigvee_{k=0}^{n-1}T^{-k}\alpha\right)$. Show by direct computation that for the Bernoulli shift $B(p, 1-p)$ with the generating partition $\alpha = \{[0], [1]\}$, the sequence $H_n/n$ is constant (i.e., $H_n = nH(\alpha)$ for all $n$). Why does this happen? (Hint: independence.)
\end{exercise}


\begin{exercise}
(Entropy of a Markov chain). Consider the Markov shift on $\{0, 1\}^{\mathbb{N}}$ with transition matrix $P = \bigl(\begin{smallmatrix} 1/2 & 1/2 \\ 1/3 & 2/3 \end{smallmatrix}\bigr)$ and stationary distribution $\pi = (2/5, 3/5)$. Show that the measure-theoretic entropy of the shift with respect to the Markov measure is

\[
h = -\sum_i \pi_i \sum_j P_{ij} \log P_{ij}.
\]

Compute this value numerically. Is it larger or smaller than the entropy of the Bernoulli shift $B(2/5, 3/5)$? Explain why.
\end{exercise}


\begin{exercise}
Let $T \colon \mathbb{T}^2 \to \mathbb{T}^2$ be the toral automorphism given by the matrix $A = \bigl(\begin{smallmatrix} 3 & 2 \\ 1 & 1\end{smallmatrix}\bigr)$. Compute the Lyapunov exponents and, using Pesin's formula, determine $h_\lambda(T)$.
\end{exercise}


\begin{exercise}
(Topological entropy of a subshift). Let $\Sigma \subset \{0, 1\}^{\mathbb{Z}}$ be the subshift of finite type defined by forbidding the word $11$ (i.e., no two consecutive $1$'s). Show that the number of admissible words of length $n$ is the Fibonacci number $F_{n+2}$, and conclude that

\[
h_{\mathrm{top}}(\sigma|_\Sigma) = \log\!\left(\frac{1 + \sqrt{5}}{2}\right).
\]
\end{exercise}


\begin{exercise}
(Variational principle application). The full shift on $\{0, 1\}^{\mathbb{Z}}$ has topological entropy $\log 2$. The Bernoulli measure $B(p, 1-p)$ has metric entropy $-p\log p - (1-p)\log(1-p)$. Verify that this is maximized over $p \in [0,1]$ at $p = 1/2$ with value $\log 2$, confirming the variational principle in this case.
\end{exercise}


\begin{exercise}
Prove that if $T$ is an isometry of a compact metric space $(X, d)$, then $h_{\mathrm{top}}(T) = 0$. (Hint: $d_n = d$ for all $n$.)
\end{exercise}


\begin{exercise}
(Entropy and factors). Let $\pi \colon (X, \mu, T) \to (Y, \nu, S)$ be a factor map (i.e., $\pi$ is measure-preserving and $\pi \circ T = S \circ \pi$ a.e.). Show that $h_\nu(S) \leq h_\mu(T)$. (Hint: if $\beta$ is a partition of $Y$, then $\pi^{-1}\beta$ is a partition of $X$ with $H_\mu(\pi^{-1}\beta) = H_\nu(\beta)$.)


\bigskip

\end{exercise}


\section*{Recommended Reading}

\begin{itemize}[nosep]
  \item \textbf{Walters, P.} \textit{An Introduction to Ergodic Theory}. Springer, 1982. --- Chapter 4 (entropy) and Chapter 8 (topological entropy and the variational principle) are the standard graduate-level reference.

  \item \textbf{Petersen, K.} \textit{Ergodic Theory}. Cambridge University Press, 1983. --- Chapters 4--6 cover entropy, generators, and Ornstein theory with careful proofs.

  \item \textbf{Downarowicz, T.} \textit{Entropy in Dynamical Systems}. Cambridge University Press, 2011. --- A modern and comprehensive treatment of both measure-theoretic and topological entropy.

  \item \textbf{Katok, A. and Hasselblatt, B.} \textit{Introduction to the Modern Theory of Dynamical Systems}. Cambridge University Press, 1995. --- Chapter 4 covers entropy; the supplements discuss Pesin theory and Lyapunov exponents.

  \item \textbf{Shannon, C. E.} "A Mathematical Theory of Communication." \textit{Bell System Technical Journal}, 27(3):379--423, 1948. --- The founding paper of information theory.

  \item \textbf{Kolmogorov, A. N.} "A New Metric Invariant of Transitive Dynamical Systems and Automorphisms of Lebesgue Spaces." \textit{Doklady Akademii Nauk SSSR}, 119(5):861--864, 1958. --- Where measure-theoretic entropy was introduced.

  \item \textbf{Ornstein, D. S.} "Bernoulli Shifts with the Same Entropy are Isomorphic." \textit{Advances in Mathematics}, 4(3):337--352, 1970. --- The landmark classification theorem.

  \item \textbf{Ornstein, D. S.} \textit{Ergodic Theory, Randomness, and Dynamical Systems}. Yale University Press, 1974. --- Ornstein's own exposition of his theory.

  \item \textbf{Brin, M. and Stuck, G.} \textit{Introduction to Dynamical Systems}. Cambridge University Press, 2002. --- An accessible introduction, covering entropy in Chapters 4--5 at a level suitable for undergraduates.


\bigskip

\end{itemize}

\textit{In the next chapter, we will explore how the ideas of ergodic theory and dynamical systems connect to reservoir computing, where the dynamics of a high-dimensional system are harnessed for computation.}
